%%--------------------------------------------------------------------
%1--------------------------------------------------------------------
\vspace*{\fill}
\begin{glyph}{Roof (屋)}{roof}\label{glyph:roof}
Roof.\\\hline
When this kanji is the final character in a compound, it will often mean a ``shop'' when referring to a building, or a ``dealer'' when associated with a person.
\end{glyph}

\begin{comps}
\comprtwocone & 尸 + 至 \\\hline
\comprthreecone & Corpse/Body + Arrive \\\hline
\end{comps}

\begin{remglyph}
The \hooglig{roof's} \remcom{body} \remcom{arrived}.\\\hline
\end{remglyph}

\begin{prons}
\pronsrtwocone & \textbf{オク (OKU)} & \textbf{や (ya)}\\\hline
\rye{3}{や is by far the more common reading, and is invariably used when this kanji appears in the final position of a compound.  There is often a mixture of on and kun-yomi in such words (as in the final 2 examples below).}
\pronsrthreecone & --- & --- \\\hline
\end{prons}

\begin{remprons}
\hooglig{Roof's} \comon{okurigana} is \comkun{ya}.\\\hline
\end{remprons}

%{\middel}
% &  \mbox{()}\\\hline
\begin{somme}
屋上 & roof + upper = \textit{rooftop} & オク{\middel}ジョウ \mbox{(OKU{\middel}J\={O})}\\\hline
本屋 & main (book) + roof = \textit{bookstore} & ホン{\middel}や \mbox{(HON{\middel}ya)}\\\hline
\notyet{料理屋} & fee + reason + roof = \textit{restaurant} & リョウ{\middel}リ{\middel}や \mbox{(RY\={O}{\middel}RI{\middel}ya)}\\\hline
\end{somme}
\end{multicols}
\vhrulefill{5pt}
%%--------------------------------------------------------------------
%2--------------------------------------------------------------------
\begin{multicols}{2}
\begin{glyph}{Each (各)}{each}\label{glyph:each}
Each.
\end{glyph}

\begin{comps}
\comprtwocone & 夂 + 口 \\\hline
\comprthreecone & Winter + Mouth\\\hline
\end{comps}

\begin{remglyph}
\hooglig{Each} \remcom{vampire's} \remcom{winter}.\\\hline
\end{remglyph}

\begin{prons}
\pronsrtwocone & \textbf{カク (KAKU)} & --- \\\hline
\pronsrthreecone & --- & おのおの　(on\={o}no) \\\hline
\end{prons}

\begin{remprons}
\hooglig{Each} \comon{kakuro} game ends with \ucomkun{``oh no, oh no!''}\\\hline
\end{remprons}

%{\middel}
% &  \mbox{()}\\\hline
\begin{somme}
各人 & each + person = \textit{each person} & カク{\middel}ジン \mbox{(KAKU{\middel}JIN)}\\\hline
\notyet{各地} & each + ground = \textit{each area} & カク{\middel}チ \mbox{(KAKU{\middel}CHI)}\\\hline
\end{somme}
\end{multicols}
\vspace*{\fill}
\pagebreak
%%--------------------------------------------------------------------
%3--------------------------------------------------------------------
\vspace*{\fill}
\begin{multicols}{2}
\begin{glyph}{Form (形)}{form}\label{glyph:form}
Form, Shape, Appearance.
\end{glyph}

\begin{comps}
\comprtwocone & 一 + 廾 + 彡 \\\hline
\comprthreecone & One + 2 hands + Hair/Feather/Bristle \\\hline
\end{comps}

\begin{remglyph}
\hooglig{Form} \remcom{two hands} out of that \remcom{one} \remcom{feather}.\\\hline
\end{remglyph}

\begin{prons}
\pronsrtwocone & \textbf{ケイ (KEI)} & \textbf{かたち、かた (katachi, kata)}\\\hline
\rye{3}{かたち is the choice when this kanji appears on its own; in such a case it means ``form'' or ``shape''.  In compound, look for ケイ in the first position (with a notable exception being the second example below), and a split between ケイ and かた after that, keeping in mind that かた is almost always voiced in such instances.}
\pronsrthreecone & ギョウ (GY\={O}) & --- \\\hline
\rye{3}{This is a very uncommon reading that you are only likely to meet in the word for `doll', 人形 (NIN{\middel}GY\={O}/ニン{\middel}ギョウ) from the kanji for `person' and `form'.}
\end{prons}

\begin{remprons}
The \hooglig{form} \comkun{catalogue} is \ucomon{gyone} with \comkun{car-touching} \comon{catering}.\\\hline
\end{remprons}

%{\middel}
% &  \mbox{()}\\\hline
\begin{somme}
形 & \textit{form; shape} & かたち \mbox{(katachi)}\\\hline
形見 & form + see = \textit{memento; keepsake} & かた{\middel}み \mbox{(kata{\middel}mi)}\\\hline
手形 & hand + form = \textit{bill; draft; note} & て{\middel}がた \mbox{(te{\middel}gata)}\\\hline
\notyet{形成} & form + become = \textit{formation} & ケイ{\middel}セイ \mbox{(KEI{\middel}SEI)}\\\hline
\end{somme}
\end{multicols}
\vhrulefill{5pt}
%%--------------------------------------------------------------------
%4--------------------------------------------------------------------
\begin{multicols}{2}
\begin{glyph}{Target (的)}{target}\label{glyph:target}
Target.\\\hline
This is an important character that is used as a suffix to express such ideas as the English equivalent of ``-istic'' in word such as ``artistic''.\\\hline
If you spend some time getting to know this kanji, you will quickly notice your vocabulary growing in sophistication.
\end{glyph}

\begin{comps}
\comprtwocone & 白 + 勹 + 丶 \\\hline
\comprthreecone & White + Embracing/Wrap + Dot \\\hline
\end{comps}

\begin{remglyph}
\hooglig{Target} the \remcom{white} \remcom{embracing} \remcom{dot}.\\\hline
\end{remglyph}

\begin{prons}
\pronsrtwocone & \textbf{テキ (TEKI)} & --- \\\hline
\pronsrthreecone & --- & まと (mato) \\\hline
\end{prons}

\begin{remprons}
The \hooglig{target} was the \ucomkun{ma-to-be} wearing \comon{tekkies}.\\\hline
\end{remprons}

%{\middel}
% &  \mbox{()}\\\hline
\begin{somme}
\rye{3}{Recall from Entry \ref{glyph:literature} that ``literature'' plus ``change'' in the second example below means ``culture'', and from Entry \ref{glyph:tool} that ``tool'' plus ``body'' in the third example means ``definite; concrete''.}
目的 & eye + target = \textit{purpose; aim} & モク{\middel}テキ \mbox{(MOKU{\middel}TEKI)}\\\hline
文化的 & culture + target = \textit{cultural} & ブン{\middel}カ{\middel}テキ \mbox{(BUN{\middel}KA{\middel}TEKI)}\\\hline
具体的 & definite + target = \textit{definite; specific; concrete} & グ{\middel}タイ{\middel}テキ \mbox{(GU{\middel}TAI{\middel}TEKI)}\\\hline
\end{somme}
\end{multicols}
\vspace*{\fill}
\pagebreak
%%--------------------------------------------------------------------
%5--------------------------------------------------------------------
\vspace*{\fill}
\begin{multicols}{2}
\begin{glyph}{Arrow (矢)}{arrow}\label{glyph:arrow}
Arrow.\\\hline
As these are the same two components used to make the kanji for ``Lose'' (Entry \ref{glyph:lose}), ensure that you are able to tell them apart.
\end{glyph}

\begin{comps}
\comprtwocone & 矢 \\\hline
\comprthreecone & Arrow \\\hline
\end{comps}

\begin{remglyph}
Part of the Kanji radicals (\#{111}).\\\hline
\end{remglyph}

\begin{prons}
\pronsrtwocone & --- & \textbf{や (ya)}\\\hline
\pronsrthreecone & シ　(SHI) & --- \\\hline
\end{prons}

\begin{remprons}
\hooglig{Arrows} \ucomon{shielded} away from the \comkun{young}.\\\hline
\end{remprons}

%{\middel}
% &  \mbox{()}\\\hline
\begin{somme}
矢 & \textit{arrow} & や \mbox{(ya)}\\\hline
弓矢 & bow + arrow = \textit{bow and arrow} & ゆみ{\middel}や \mbox{(yumi{\middel}ya)}\\\hline
\end{somme}
\end{multicols}
\vhrulefill{5pt}
%%--------------------------------------------------------------------
%6--------------------------------------------------------------------
\begin{multicols}{2}
\begin{glyph}{World (界)}{world}\label{glyph:world}
World.\\\hline
Circles (in the sense of ``financial circles'' or ``scientific circles'', etc.).
\end{glyph}

\begin{comps}
\comprtwocone & 田 + 介 \\\hline
\comprthreecone & Rice Field + Intervene (\ref{glyph:intervene}) \\\hline
\end{comps}

\begin{remglyph}
The \hooglig{world} is \remcom{intervening} with our \remcom{rice} fields.\\\hline
\end{remglyph}

\begin{prons}
\pronsrtwocone & \textbf{カイ (KAI)} & ---\\\hline
\pronsrthreecone & --- & --- \\\hline
\end{prons}

\begin{remprons}
A \hooglig{world} for \comon{kite} fanatics.\\\hline
\end{remprons}

%{\middel}
% &  \mbox{()}\\\hline
\begin{somme}
世界 & society + world = \textit{the world} & セ{\middel}カイ \mbox{(SE{\middel}KAI)}\\\hline
\notyet{魚介} & business + world = \textit{the business world} & ギョ{\middel}カイ \mbox{(GYO{\middel}KAI)}\\\hline
\notyet{自然界} & self + nature + world = \textit{the world of nature} & シ{\middel}ゼン{\middel}カイ \mbox{(SHI{\middel}ZEN{\middel}KAI)}\\\hline
\end{somme}
\end{multicols}
\vspace*{\fill}
\pagebreak
%%--------------------------------------------------------------------
%7--------------------------------------------------------------------
\vspace*{\fill}
\begin{multicols}{2}
\begin{glyph}{Against (反)}{against}\label{glyph:against}
Against.\\\hline
Anti-
\end{glyph}

\begin{comps}
\comprtwocone & 厂 + 又 \\\hline
\comprthreecone & Cliff + Again\ \\\hline
\end{comps}

\begin{remglyph}
\hooglig{Against} \remcom{cliffs} \remcom{again}.\\\hline
\end{remglyph}

\begin{prons}
\pronsrtwocone & \textbf{ハン (HAN)} & --- \\\hline
\pronsrthreecone & タン、ホン (TAN, HON) & そ (so) \\\hline
\end{prons}

\begin{remprons}
\hooglig{Against} \ucomkun{softies} \ucomon{haunted} by \ucomon{tons} of \comon{hunches}.\\\hline
\end{remprons}

%{\middel}
% &  \mbox{()}\\\hline
\begin{somme}
反転 & against + roll = \textit{roll over; reverse} & ハン{\middel}テン \mbox{(HAN{\middel}TEN)}\\\hline
\notyet{反動} & against + move = \textit{reaction; recoil} & ハン{\middel}ドウ \mbox{(HAN{\middel}D\={O})}\\\hline
\notyet{反対} & against + versus = \textit{opposition; resistance} & ハン{\middel}タイ \mbox{(HAN{\middel}TAI)}\\\hline
\end{somme}
\end{multicols}
\vhrulefill{5pt}
%%--------------------------------------------------------------------
%8--------------------------------------------------------------------
\begin{multicols}{2}
\begin{glyph}{Good (良)}{good}\label{glyph:good}
Good.
\end{glyph}

\begin{comps}
\comprtwocone & 艮 \\\hline
\comprthreecone & Boundary \\\hline
\end{comps}

\begin{remglyph}
Part of the Kanji radicals (\#{138}).\\\hline
\end{remglyph}

\begin{prons}
\pronsrtwocone & \textbf{リョウ (RY\={O})} & \textbf{よ (yo)}\\\hline
\pronsrthreecone & --- & --- \\\hline
\end{prons}

\begin{remprons}
A \hooglig{good} \comkun{yomp} with \comon{Ryobi} gear.\\\hline
\end{remprons}

%{\middel}
% &  \mbox{()}\\\hline
\begin{somme}
良い & \textit{good} & よ{\middel}い \mbox{(yo{\middel}i)}\\\hline
良心 & good + heart = \textit{conscience} & リョウ{\middel}シン \mbox{(RY\={O}{\middel}SHIN)}\\\hline
良好 & good + like = \textit{satisfactory} & リョウ{\middel}コ \mbox{(RY\={O}{\middel}KO)}\\\hline
\end{somme}
\end{multicols}
\vspace*{\fill}
\pagebreak
%%--------------------------------------------------------------------
%9--------------------------------------------------------------------
\vspace*{\fill}
\begin{multicols}{2}
\begin{glyph}{Distribute (配)}{distribute}\label{glyph:distribute}
Distribute.
\end{glyph}

\begin{comps}
\comprtwocone & 酉　+ 己 \\\hline
\comprthreecone & Wine/Jar + Oneself \\\hline
\end{comps}

\begin{remglyph}
\hooglig{Distribute} some \remcom{wine} to \remcom{oneself}.\\\hline
\end{remglyph}

\begin{prons}
\pronsrtwocone & \textbf{ハイ (HAI)} & \textbf{くば (kuba)}\\\hline
\pronsrthreecone & --- & --- \\\hline
\end{prons}

\begin{remprons}
\hooglig{Distribute} \comon{hiking} gear to \comkun{Cuba}.\\\hline
\end{remprons}

%{\middel}
% &  \mbox{()}\\\hline
\begin{somme}
配る & \textit{distribute} & くば{\middel}る \mbox{(kuba{\middel}ru)}\\\hline
配水 & distribute + water = \textit{water supply} & ハイ{\middel}スイ \mbox{(HAI{\middel}SUI)}\\\hline
心配 & heart + distribute = \textit{worry; anxiety} & シン{\middel}パイ \mbox{(SHIN{\middel}PAI)}\\\hline
\end{somme}
\end{multicols}
\vhrulefill{5pt}
%%--------------------------------------------------------------------
%10-------------------------------------------------------------------
\begin{multicols}{2}
\begin{glyph}{Far (遠)}{far}\label{glyph:far}
Far.
\end{glyph}

\begin{comps}
\comprtwocone & 土 + 辶 + 口 + 衣 \\\hline
\comprthreecone & Earth + Walking/Road + Mouth + Clothing \\\hline
\end{comps}

\begin{remglyph}
\hooglig{Far} from \remcom{Earth}, \remcom{vampires} are \remcom{walking} to get \remcom{clothing}.\\\hline
\end{remglyph}

\begin{prons}
\pronsrtwocone & \textbf{エン (EN)} & \textbf{とお (t\={o})}\\\hline
\pronsrthreecone & オン (ON) & --- \\\hline
\end{prons}

\begin{remprons}
The \hooglig{far} \comon{entrance} will \comkun{torment} you \ucomon{online}.\\\hline
\end{remprons}

%{\middel}
% &  \mbox{()}\\\hline
\begin{somme}
遠い & \textit{far; distant} & とお{\middel}い \mbox{(t\={o}{\middel}i)}\\\hline
遠足 & far + leg = \textit{excursion} & エン{\middel}ソク \mbox{(EN{\middel}SOKU)}\\\hline
\notyet{永遠} & eternal + far = \textit{eternity} & エイ{\middel}エン \mbox{(EI{\middel}EN)}\\\hline
\end{somme}
\end{multicols}
\vspace*{\fill}
\pagebreak
%%--------------------------------------------------------------------
%11-------------------------------------------------------------------
\vspace*{\fill}
\begin{multicols}{2}
\begin{glyph}{Forget (忘)}{forget}\label{glyph:forget}
Forget.
\end{glyph}

\begin{comps}
\comprtwocone & 亡 + 心 \\\hline
\comprthreecone & Dead/Gone + Heart \\\hline
\end{comps}

\begin{remglyph}
\hooglig{Forget} that your \remcom{heart} shall \remcom{perish}.\\\hline
\end{remglyph}

\begin{prons}
\pronsrtwocone & \textbf{ボウ (B\={O})} & \textbf{わす (wasu)}\\\hline
\pronsrthreecone & --- & --- \\\hline
\end{prons}

\begin{remprons}
\hooglig{Forget} \comon{boring} \comkun{wassups}.\\\hline
\end{remprons}

%{\middel}
% &  \mbox{()}\\\hline
\begin{somme}
忘れる & \textit{forget} & わす{\middel}れる \mbox{(wasu{\middel}reru)}\\\hline
忘れ物 & forget + thing = \textit{something left behind} & わす{\middel}れ　もの \mbox{(wasu{\middel}re mono)}\\\hline
忘年会 & forget + year + meet = \textit{year-end party} & ボウ{\middel}ネン{\middel}カイ \mbox{(B\={O}{\middel}NEN{\middel}KAI)}\\\hline
\end{somme}
\end{multicols}
\vhrulefill{5pt}
%%--------------------------------------------------------------------
%12-------------------------------------------------------------------
\begin{multicols}{2}
\begin{glyph}{Perfect (完)}{perfect}\label{glyph:perfect}
Perfect, Complete.
\end{glyph}

\begin{comps}
\comprtwocone & 宀 + 儿 + 二 \\\hline
\comprthreecone & Roof/Cover + Man/Person + Two \\\hline
\end{comps}

\begin{remglyph}
A \hooglig{perfect} \remcom{roof} for \remcom{2} \remcom{people}.\\\hline
\end{remglyph}

\begin{prons}
\pronsrtwocone & \textbf{カン (KAN)} & --- \\\hline
\pronsrthreecone & --- & --- \\\hline
\end{prons}

\begin{remprons}
\hooglig{Perfect} \comon{cunning}.\\\hline
\end{remprons}

%{\middel}
% &  \mbox{()}\\\hline
\begin{somme}
完全 & perfect + complete = \textit{perfection; completeness} & カン{\middel}ゼン \mbox{(KAN{\middel}ZEN)}\\\hline
\notyet{完成} & perfect + become = \textit{completion; accomplishment} & カン{\middel}セイ \mbox{(KAN{\middel}SEI)}\\\hline
\notyet{完結} & perfect + bind = \textit{completion; conclusion} & カン{\middel}ケツ \mbox{(KAN{\middel}KETSU)}\\\hline
\end{somme}
\end{multicols}
\vspace*{\fill}
\pagebreak
%%--------------------------------------------------------------------
%13-------------------------------------------------------------------
\vspace*{\fill}
\begin{multicols}{2}
\begin{glyph}{Besides (且)}{besides}\label{glyph:besides}
Besides, Moreover.
\end{glyph}

\begin{comps}
\comprtwocone & 目 + 一 \\\hline
\comprthreecone & Eye + One \\\hline
\end{comps}

\begin{remglyph}
\hooglig{Besides} the \remcom{one} \remcom{eye}.\\\hline
\end{remglyph}

\begin{prons}
\pronsrtwocone & --- & \textbf{か (ka)}\\\hline
\pronsrthreecone & ショ (SHO) & --- \\\hline
\end{prons}

\begin{remprons}
\hooglig{Besides} the \ucomon{shopping} \comkun{customer}.\\\hline
\end{remprons}

%{\middel}
% &  \mbox{()}\\\hline
\begin{somme}
且つ & \textit{besides} & か{\middel}つ \mbox{(ka{\middel}tsu)}\\\hline
\end{somme}
\end{multicols}
\vhrulefill{5pt}
%%--------------------------------------------------------------------
%14-------------------------------------------------------------------
\begin{multicols}{2}
\begin{glyph}{Metropolis (都)}{metropolis}\label{glyph:metropolis}
Metropolis.  Large City.
\end{glyph}

\begin{comps}
\comprtwocone & 老耂 + 阝邑 + 日 \\\hline
\comprthreecone & Old + City + Sun \\\hline
\end{comps}

\begin{remglyph}
The \hooglig{metropolis} was a \remcom{sunny} \remcom{old} \remcom{city}.\\\hline
\end{remglyph}

\begin{prons}
\pronsrtwocone & \textbf{ト (TO)} & \textbf{みやこ (miyako)}\\\hline
\pronsrthreecone & ツ (TSU) & --- \\\hline
\rye{3}{You will only meet this on-yomi in two compounds: the first being 都合 (ツ{\middel}ゴウ/TSU{\middel}G\={O}) \textit{``circumstances''}, \textit{``arrangements''}, composed of the kanji ``metropolis'' and ``match''.\\\hline
The second compound is 都度 (ツ{\middel}ド/TSU{\middel}DO) \textit{``each time''}/\textit{``whenever''} from the kanji ``metropolis'' and ``degree''.}
\end{prons}

\begin{remprons}
The \hooglig{metropolis} \comkun{Mia completed} was \comon{toppled} by the \ucomon{tsunami}.\\\hline
\end{remprons}

%{\middel}
% &  \mbox{()}\\\hline
\begin{somme}
都 & \textit{metropolis; capital} & みやこ \mbox{(miyako)}\\\hline
京都 & capital + metropolis = \textit{Kyoto} & キョウ{\middel}ト \mbox{(KY\={O}{\middel}TO)}\\\hline
首都 & neck + metropolis = \textit{capital city} & シュ{\middel}ト \mbox{(SHU{\middel}TO)}\\\hline
\end{somme}
\end{multicols}
\vspace*{\fill}
\pagebreak
%%--------------------------------------------------------------------
%15-------------------------------------------------------------------
\vspace*{\fill}
\begin{multicols}{2}
\begin{glyph}{Descend (降)}{descend}\label{glyph:descend}
Descend, fall, surrender.\\\hline
This is the character used when describing the `falling' of rain or snow, etc.
\end{glyph}

\begin{comps}
\comprtwocone & 十 + 阝阜 + 夂 \\\hline
\comprthreecone & Ten + Hill/Dam/Mound + Winter \\\hline
\end{comps}

\begin{remglyph}
The \hooglig{descending} \remcom{scarecrow} was \remcom{wintering} in the \remcom{mound}.\\\hline
\end{remglyph}

\begin{prons}
\pronsrtwocone & \textbf{コウ (K\={O})} & \textbf{ふ (fu)}\\\hline
\pronsrthreecone & --- & お (o) \\\hline
\rye{3}{This reading forms an intransitive/transitive verb pair (降りる/降ろす) that has the same pronunciation as one we met earlier with the kanji 下 (Entry 31): 下りる/下ろす.  The characters are sometimes used interchangeably, although you are more likely to see then written as 下りる and 下ろす.}
\end{prons}

\begin{remprons}
\hooglig{Descending} scores \comon{caused} the \comkun{foosball} team to be \ucomkun{obligated} to do something.\\\hline
\end{remprons}

%{\middel}
% &  \mbox{()}\\\hline
\begin{somme}
降る & \textit{falling (rain/snow)}　& ふ{\middel}る \mbox{(fu{\middel}ru)}\\\hline
降下 & descend + lower = \textit{descent; landing}　& コウ{\middel}カ \mbox{(K\={O}{\middel}KA)}\\\hline
降水 & descend + water = \textit{precipitation}　& コウ{\middel}スイ \mbox{(K\={O}{\middel}SUI)}\\\hline
\end{somme}
\end{multicols}
\vhrulefill{5pt}
%%--------------------------------------------------------------------
%16-------------------------------------------------------------------
\begin{multicols}{2}
\begin{glyph}{Marriage (婚)}{marriage}\label{glyph:marriage}
Marriage.
\end{glyph}

\begin{comps}
\comprtwocone & 日 + 女 + 氏 \\\hline
\comprthreecone & Sun + Woman/Female + Clan \\\hline
\end{comps}

\begin{remglyph}
\hooglig{Marrying} a \remcom{woman} into a \remcom{clan} is as clear as \remcom{day}.\\\hline
\end{remglyph}

\begin{prons}
\pronsrtwocone & \textbf{コン (KON)} & --- \\\hline
\pronsrthreecone & --- & --- \\\hline
\end{prons}

\begin{remprons}
\hooglig{Marriage} \comon{con}.\\\hline
\end{remprons}

%{\middel}
% &  \mbox{()}\\\hline
\begin{somme}
離婚 & separage + marriage = \textit{divorce} & リ{\middel}コン \mbox{(RI{\middel}KON)}\\\hline
\notyet{結婚} & bind + marriage = \textit{marriage} & ケッ{\middel}コン \mbox{(KEK{\middel}KON)}\\\hline
\notyet{婚期} & marriage + period = \textit{marriageable age} & コン{\middel}キ \mbox{(KON{\middel}KI)}\\\hline
\end{somme}
\end{multicols}
\vspace*{\fill}
\pagebreak
%%--------------------------------------------------------------------
%17-------------------------------------------------------------------
\vspace*{\fill}
\begin{multicols}{2}
\begin{glyph}{Festival (祭)}{festival}\label{glyph:festival}
Festival.\\\hline
This kanji can also have some religious overtones in the senses of ritual and worship.\\\hline
Note that the first stroke of `again' is clipped off in some fonts so that it does not pass through the second.
\end{glyph}

\begin{comps}
\comprtwocone & 月 + 又 + 示 \\\hline
\comprthreecone & Moon + Again + Show \\\hline
\end{comps}

\begin{remglyph}
\hooglig{Festivals} \remcom{again} on \remcom{show} this \remcom{month}.\\\hline
\end{remglyph}

\begin{prons}
\pronsrtwocone & \textbf{サイ (SAI)} & \textbf{まつ (matsu)}\\\hline
\pronsrthreecone & --- & --- \\\hline
\end{prons}

\begin{remprons}
A \hooglig{festival} of \comon{psycho} \comkun{mutts}.\\\hline
\end{remprons}

%{\middel}
% &  \mbox{()}\\\hline
\begin{somme}
祭り & \textit{festival} & まつ{\middel}り \mbox{(matsu{\middel}ri)}\\\hline
祭日 & festival + sun (day) = \textit{festival day} & サイ{\middel}ジツ \mbox{(SAI{\middel}JITSU)}\\\hline
百年祭 & hundred + year + festival = \textit{centennial} & ヒャク{\middel}ネン{\middel}サイ \mbox{(HYAKU{\middel}NEN{\middel}SAI)}\\\hline
\end{somme}
\end{multicols}
\vhrulefill{5pt}
%%--------------------------------------------------------------------
%18-------------------------------------------------------------------
\begin{multicols}{2}
\begin{glyph}{Finger (指)}{finger}\label{glyph:finger}
Finger, along with the senses of indicating and pointing.\\\hline
Note that this kanji never appears as a component in other characters.
\end{glyph}

\begin{comps}
\comprtwocone & 手扌才 + 旨 \\\hline
\comprthreecone & Hand/Actions + Gist (\ref{glyph:gist}) \\\hline
\end{comps}

\begin{remglyph}
Your \hooglig{finger} shows the \remcom{gist} of your \remcom{hand actions}.\\\hline
\end{remglyph}

\begin{prons}
\pronsrtwocone & \textbf{シ (SHI)} & \textbf{ゆび、さ (yubi, sa)}\\\hline
\rye{3}{As the second compound below shows, さ can become voiced outside of the first position.}
\pronsrthreecone & --- & --- \\\hline
\end{prons}

\begin{remprons}
His \hooglig{finger} indicated to a \comkun{summary} on the \comon{shield} which stated: "\comkun{you before} me"\\\hline
\end{remprons}

%{\middel}
% &  \mbox{()}\\\hline
\begin{somme}
指 & \textit{finger} & ゆび \mbox{(yubi)}\\\hline
目指す & eye + finger = \textit{aim at} & め　ざす \mbox{(me za{\middel}su)}\\\hline
指名 & finger + name = \textit{nomination; designation} & シ{\middel}メイ \mbox{(SHI{\middel}MEI)}\\\hline
\end{somme}
\end{multicols}
\vspace*{\fill}
\pagebreak
%%--------------------------------------------------------------------
%19-------------------------------------------------------------------
\vspace*{\fill}
\begin{multicols}{2}
\begin{glyph}{Beforehand (予)}{beforehand}\label{glyph:beforehand}
Beforehand, previously.
\end{glyph}

\begin{comps}
\comprtwocone & 亅 + 一 + マ \\\hline
\comprthreecone & Hook + One + MA katakana \\\hline
\end{comps}

\begin{remglyph}
\hooglig{Beforehand} \remcom{mum} took \remcom{one} \remcom{hook}.\\\hline
\end{remglyph}

\begin{prons}
\pronsrtwocone & \textbf{ヨ (YO)} & --- \\\hline
\pronsrthreecone & --- & --- \\\hline
\end{prons}

\begin{remprons}
\hooglig{Previous} \comon{yob}.\\\hline
\end{remprons}

%{\middel}
% &  \mbox{()}\\\hline
\begin{somme}
予言 & beforehand + say = \textit{prediction} & ヨ{\middel}ゲン \mbox{(YO{\middel}GEN)}\\\hline
予定 & beforehand + fixed = \textit{prearrangement; plan} & ヨ{\middel}テイ \mbox{(YO{\middel}TEI)}\\\hline
\notyet{予期} & beforehand + period = \textit{anticipation; expectation} & ヨ{\middel}キ \mbox{(YO{\middel}KI)}\\\hline
\end{somme}
\end{multicols}
\vhrulefill{5pt}
%%--------------------------------------------------------------------
%20-------------------------------------------------------------------
\begin{multicols}{2}
\begin{glyph}{Active (活)}{active}\label{glyph:active}
Active, Lively.
\end{glyph}

\begin{comps}
\comprtwocone & 舌 + 水氵氺 \\\hline
\comprthreecone & Tongue + Water \\\hline
\end{comps}

\begin{remglyph}
\hooglig{Actively} the \remcom{tongue} was taking \remcom{water} in.\\\hline
\end{remglyph}

\begin{prons}
\pronsrtwocone & \textbf{カツ (KATSU)} & --- \\\hline
\pronsrthreecone & --- & --- \\\hline
\end{prons}

\begin{remprons}
\hooglig{Active}, lively \comon{cuts}.\\\hline
\end{remprons}

%{\middel}
% &  \mbox{()}\\\hline
\begin{somme}
生活 & life + active = \textit{life; livelihood} & セイ{\middel}カツ \mbox{(SEI{\middel}KATSU)}\\\hline
活気 & active + spirit = \textit{liveliness; vigour} & カッ{\middel}キ \mbox{(KAK{\middel}KI)}\\\hline
\notyet{活動} & active + move = \textit{activity; action} & カツ{\middel}ドウ \mbox{(KATSU{\middel}D\={O})}\\\hline
\end{somme}
\end{multicols}
\vspace*{\fill}
\pagebreak
%%--------------------------------------------------------------------
%21-------------------------------------------------------------------
\vspace*{\fill}
\begin{multicols}{2}
\begin{glyph}{Harmony (和)}{harmony}\label{glyph:harmony}
Harmony, peace.\\\hline
This kanji is also used as an abbreviation for Japan in much the same way 日 is (the final compound offers an example of this usage).
\end{glyph}

\begin{comps}
\comprtwocone & 口 + 禾 \\\hline
\comprthreecone & Mouth/Opening/Sounding + Grain (Ear) \\\hline
\end{comps}

\begin{remglyph}
\hooglig{Harmonious} \remcom{sounding} of grinding \remcom{grain}.\\\hline
\end{remglyph}

\begin{prons}
\pronsrtwocone & \textbf{ワ (WA)} & --- \\\hline
\pronsrthreecone & オ　(O) & やわ、なご (yawa, nago) \\\hline
\end{prons}

\begin{remprons}
\hooglig{Harmonious} \comon{washing} of \comkun{Nago} \ucomon{objects} with \comkun{ya water}.\\\hline
\end{remprons}

%{\middel}
% &  \mbox{()}\\\hline
\begin{somme}
和 & \textit{harmony} & ワ \mbox{(WA)}\\\hline
\notyet{平和} & flat + harmony = \textit{peace} & ヘイ{\middel}ワ \mbox{(HEI{\middel}WA)}\\\hline
\notyet{和食} & harmony (Japan) + eat = \textit{Japanese food} & ワ{\middel}ショク \mbox{(WA{\middel}SHOKU)}\\\hline
\end{somme}

\begin{irrr}
	大和 & large + harmony = \textit{Ancient Japan} & やまと \mbox{(yamato)}\\\hline
	日和 & sun + harmony = \textit{fine weather} & ひ{\middel}より \mbox{(hi{\middel}yori)}\\\hline
\end{irrr}
\end{multicols}
\vhrulefill{5pt}
%%--------------------------------------------------------------------
%22-------------------------------------------------------------------
\begin{multicols}{2}
\begin{glyph}{Via (経)}{via}\label{glyph:via}
Via, in the sense of things passing or elapsing.\\\hline
A secondary meaning related to Buddhist sutras make use of the less common on-yomi.
\end{glyph}

\begin{comps}
\comprtwocone & 糸 + 又 + 土 \\\hline
\comprthreecone & Thread + Again + Earth \\\hline
\end{comps}

\begin{remglyph}
\hooglig{Via} \remcom{threading} the \remcom{Earth} \remcom{again}.\\\hline
\end{remglyph}

\begin{prons}
\pronsrtwocone & \textbf{ケイ (KEI)} & \textbf{へ (he)}\\\hline
\pronsrthreecone & キョウ (KY\={O}) & --- \\\hline
\end{prons}

\begin{remprons}
\hooglig{Via} \comkun{headcount} we \comon{catered} in \ucomon{Kyoto}.\\\hline
\end{remprons}

%{\middel}
% &  \mbox{()}\\\hline
\begin{somme}
経る & \textit{elapse; pass through} & へ{\middel}る \mbox{(he{\middel}ru)}\\\hline
経由 & via + cause = \textit{via; by way of} & ケイ{\middel}ユ \mbox{(KEI{\middel}YU)}\\\hline
\notyet{経験} & via + test = \textit{experience} & ケイ{\middel}ケン \mbox{(KEI{\middel}KEN)}\\\hline
\end{somme}
\end{multicols}
\vspace*{\fill}
\pagebreak
%%--------------------------------------------------------------------
%23-------------------------------------------------------------------
\vspace*{\fill}
\begin{multicols}{2}
\begin{glyph}{Head (頭)}{head}\label{glyph:head}
Head.
\end{glyph}

\begin{comps}
\comprtwocone & 豆 + 頁 \\\hline
\comprthreecone & Bean/A kind of a Vessel + Big Shell \\\hline
\end{comps}

\begin{remglyph}
\hooglig{Head} is \remcom{a kind of a vessel} like a \remcom{big shell}.\\\hline
\end{remglyph}

\begin{prons}
\pronsrtwocone & \textbf{トウ (T\={O})} & \textbf{あたま (atama)}\\\hline
\pronsrthreecone & ズ、ト (ZU, TO) & かしら (kashira) \\\hline
\rye{3}{You are most likely to encounter ズ in only a pair of words: 頭上 (ズ{\middel}ジョウ/ZU{\middel}J\={O}) \textit{``overhead''}, from the kanji `head' and `upper', and the final sample compound of Entry 461.}
\end{prons}

\begin{remprons}
\hooglig{Head} \ucomon{tops} of \comkun{atamasco} lillies \comon{tortured} by the \ucomkun{cashier} into \ucomon{Zucchini} strips.\\\hline
\end{remprons}

%{\middel}
% &  \mbox{()}\\\hline
\begin{somme}
頭 & \textit{head} & あたま \mbox{(atama)}\\\hline
先頭 & precede + head = \textit{at the lead; at the head} & セン{\middel}トウ \mbox{(SEN{\middel}T\={O})}\\\hline
\notyet{店頭} & shop + head = \textit{storefront; shopwindow} & テン{\middel}トウ \mbox{(TEN{\middel}T\={O})}\\\hline
\end{somme}
\end{multicols}
\vhrulefill{5pt}
%%--------------------------------------------------------------------
%24-------------------------------------------------------------------
\begin{multicols}{2}
\begin{glyph}{Fish (魚)}{fish}\label{glyph:fish}
Fish.
\end{glyph}

\begin{comps}
\comprtwocone & 魚 \\\hline
\comprthreecone & Fish \\\hline
\end{comps}

\begin{remglyph}
Part of the Kanji radicals (\#{195}).\\\hline
\end{remglyph}

\begin{prons}
\pronsrtwocone & \textbf{ギョ (GYO)} & \textbf{さかな (sakana)}\\\hline
\pronsrthreecone & --- & うお (uo) \\\hline
\end{prons}

\begin{remprons}
\hooglig{Fish} \comkun{suck on arteries} \ucomkun{you occupy} in \comon{Gion}.\\\hline
\end{remprons}

%{\middel}
% &  \mbox{()}\\\hline
\begin{somme}
魚 & \textit{fish} & さかな \mbox{(sakana)}\\\hline
金魚 & gold + fish = \textit{goldfish} & キン{\middel}ギョ \mbox{(KIN{\middel}GYO)}\\\hline
人魚 & person + fish = \textit{mermaid; merman} & ニン{\middel}ギョ \mbox{(NIN{\middel}GYO)}\\\hline
\end{somme}
\end{multicols}
\vspace*{\fill}
\pagebreak
%%--------------------------------------------------------------------
%25-------------------------------------------------------------------
\vspace*{\fill}
\begin{multicols}{2}
\begin{glyph}{Actual (実)}{actual}\label{glyph:actual}
This kanji has to do with the actual nature or reality of things.
\end{glyph}

\begin{comps}
\comprtwocone & 宀 + 三 + 人 \\\hline
\comprthreecone & Roof + Three + Person \\\hline
\end{comps}

\begin{remglyph}
\hooglig{Actually}, the \remcom{roof} covered \remcom{three} \remcom{people}.\\\hline
\end{remglyph}

\begin{prons}
\pronsrtwocone & \textbf{ジツ (JITSU)} & --- \\\hline
\pronsrthreecone & --- & み、みの (mi, mino) \\\hline
\end{prons}

\begin{remprons}
\hooglig{Actually} a \ucomkun{minor} \ucomkun{mirror} \comon{Ju-Jitsu}.\\\hline
\end{remprons}

%{\middel}
% &  \mbox{()}\\\hline
\begin{somme}
事実 & matter + actual = \textit{fact} & ジ{\middel}ジツ \mbox{(JI{\middel}JITSU)}\\\hline
実用 & actual + utilize = \textit{practical use; utility} & ジツ{\middel}ヨウ \mbox{(JITSU{\middel}Y\={O})}\\\hline
実行 & actual + go = \textit{put into practice; execute} & ジッ{\middel}コウ \mbox{(JIK{\middel}K\={O})}\\\hline
\end{somme}
\end{multicols}
\vhrulefill{5pt}
\vspace*{\fill}
\pagebreak
%%--------------------------------------------------------------------
